% ===================================================================
\section{Appendix C: Theory Incorporation Proposals}
\label{app:proposals}
% ===================================================================

This appendix provides the full text of all eleven theory incorporation proposals
generated during the FRFP formalization process. Each proposal is a recommended
change to the theory paper(s) to align the informal presentation with the formal
Lean development.

\begin{proposal}{P-01}
\textbf{Title}: Reduction operates on term objects, not category objects.

\noindent\textbf{Problem}: The FRFP theory paper describes reduction as operating
on ``objects in the explicit category $\Ecat$.'' In the Lean formalization, this
caused a type error: categorical objects are not terms, and reduction operates on
terms in a term calculus.

\noindent\textbf{Proposed fix}: Add a section to the theory paper distinguishing
the term calculus (the syntactic language on which reduction is defined) from the
categorical semantics (the mathematical universe in which the denotational
meanings of terms live). State explicitly that reduction arrows are term-level,
while functors are category-level.

\noindent\textbf{Impact}: Chapter 3 (reduction), Chapter 4 (semantics).
\end{proposal}

\begin{proposal}{P-02}
\textbf{Title}: The run-to-$\omega$ bijection requires an explicit construction.

\noindent\textbf{Problem}: The theory paper describes the run-to-$\omega$
correspondence informally as ``a natural correspondence.'' In Lean, we had to
construct explicit functions \texttt{runToOmega} and \texttt{omegaToRun} and
prove two round-trip properties.

\noindent\textbf{Proposed fix}: Add a formal definition of the bijection in the
theory paper, stating:
\[
\mathrm{runTo}\omega : \mathrm{Trace} \to \Omega
\qquad \text{and} \qquad
\omega\mathrm{ToRun} : \Omega \to \mathrm{Trace}
\]
with round-trip equations $\omega\mathrm{ToRun} \circ \mathrm{runTo}\omega = \mathrm{id}$
and $\mathrm{runTo}\omega \circ \omega\mathrm{ToRun} = \mathrm{id}$.

\noindent\textbf{Impact}: Chapter 4 (semantics).
\end{proposal}

\begin{proposal}{P-03}
\textbf{Title}: Credence clamping is a normative design choice, not a consequence.

\noindent\textbf{Problem}: The theory paper implies that credence values are
``naturally'' bounded in $[0,1]$. In Lean, we discovered that values outside
$[0,1]$ are type-valid without an explicit clamping operation. Clamping is a
design choice, not a mathematical consequence.

\noindent\textbf{Proposed fix}: Define the clamping function explicitly:
\[
\mathrm{clamp}(x) = \max(0, \min(1, x))
\]
and state that all credence computations apply $\mathrm{clamp}$ after each
arithmetic step. Justify the choice normatively.

\noindent\textbf{Impact}: Chapter 5 (dynamic layer).
\end{proposal}

\begin{proposal}{P-04}
\textbf{Title}: Navigation equivalence must be explicitly defined as an equivalence relation.

\noindent\textbf{Problem}: The theory paper uses $\sim_\mathrm{nav}$ without
stating reflexivity, symmetry, and transitivity. Lean requires all three proofs
explicitly.

\noindent\textbf{Proposed fix}: Add a definition block stating:
\[
t \sim_\mathrm{nav} t \quad \text{(reflexivity)},
\quad t \sim_\mathrm{nav} t' \Rightarrow t' \sim_\mathrm{nav} t \quad \text{(symmetry)},
\quad t \sim_\mathrm{nav} t'' \text{ if } t \sim_\mathrm{nav} t' \sim_\mathrm{nav} t''.
\]

\noindent\textbf{Impact}: Chapter 3 (reduction), Chapter 4 (semantics).
\end{proposal}

\begin{proposal}{P-05}
\textbf{Title}: Normal forms are unique modulo navigation equivalence, not absolutely unique.

\noindent\textbf{Problem}: The theory paper claims ``unique normal forms.'' In
Lean, we proved uniqueness only up to $\sim_\mathrm{nav}$. Two terms can have
the same observable behavior but different navigation paths; their normal forms
are navigation-equivalent but not definitionally equal.

\noindent\textbf{Proposed fix}: Correct the uniqueness claim to read ``unique
normal form up to navigation equivalence,'' and state the theorem explicitly.

\noindent\textbf{Impact}: Chapter 3 (reduction).
\end{proposal}

\begin{proposal}{P-06}
\textbf{Title}: Governance safe extension requires a step-compatibility precondition.

\noindent\textbf{Problem}: Without a compatibility precondition, a malformed step
could be appended to a well-formed trace, invalidating the extension. The
step-compatibility predicate must be defined and included in the theorem statement.

\noindent\textbf{Proposed fix}: Define \texttt{StepCompatible}(trace, step) as
the predicate that the new step's intent and authorization are coherent with the
preceding context. State the safe extension theorem explicitly.

\noindent\textbf{Impact}: Chapter 6 (governance), new theorem section.
\end{proposal}

\begin{proposal}{P-07}
\textbf{Title}: First-witness canonicity should be stated as a numbered theorem.

\noindent\textbf{Problem}: The theory paper notes informally that governance
failures have a ``first occurrence.'' This deserves formal theorem status since
it is the foundation for constructive governance monitoring algorithms.

\noindent\textbf{Proposed fix}: Add a theorem: ``For any malformed governance
trace, there exists a unique minimal index $k^*$ such that step $k^*$ is the
first governance violation.''

\noindent\textbf{Impact}: Chapter 6 (governance), new results section.
\end{proposal}

\begin{proposal}{P-08}
\textbf{Title}: Consensus impossibility requires a typed no-grounding precondition.

\noindent\textbf{Problem}: The original consensus theorem admitted unintended
models where agents could achieve consensus by vacuous means. A typed
\texttt{groundingRequired : Bool} precondition prevents this collapse.

\noindent\textbf{Proposed fix}: Add the grounding-required type flag as an explicit
hypothesis in the consensus impossibility statement.

\noindent\textbf{Impact}: Chapter 6 (collective layer).
\end{proposal}

\begin{proposal}{P-09}
\textbf{Title}: The Grothendieck construction requires a split-opfibration condition.

\noindent\textbf{Problem}: The theory paper applies the Grothendieck construction
without mentioning the split-opfibration requirement on the indexing functor. Lean
required this structure explicitly.

\noindent\textbf{Proposed fix}: Add a brief categorical prerequisite section
noting that the FRFP semantic functor $F : \Czero \to \mathbf{Cat}$ must be a
split opfibration. Cite Mac Lane (1971) Ch.~XII.

\noindent\textbf{Impact}: Chapter 4 (semantics).
\end{proposal}

\begin{proposal}{P-10}
\textbf{Title}: Phase partition well-definedness requires finiteness of pipeline terms.

\noindent\textbf{Problem}: The phase partition theorem holds only for finite
pipeline terms. Lean's type-checker requires a finiteness hypothesis; without it,
the proof does not go through.

\noindent\textbf{Proposed fix}: Add ``for any finite pipeline term $t$'' as a
hypothesis to the phase partition theorem. Note that infinite pipelines are not
currently modeled.

\noindent\textbf{Impact}: Chapter 1 (kernel), Chapter 3 (reduction).
\end{proposal}

\begin{proposal}{P-11}
\textbf{Title}: Stopping time measurability requires explicit filtration structure.

\noindent\textbf{Problem}: The theory paper cites stopping times without specifying
the filtration $(\mathcal{F}_t)_{t \geq 0}$ with respect to which measurability
is claimed. Lean's Mathlib requires the filtration as an explicit argument.

\noindent\textbf{Proposed fix}: Add a definition of the natural FRFP filtration:
$\mathcal{F}_t = \sigma(\text{all protocol events up to time } t)$, and state
measurability with respect to this filtration.

\noindent\textbf{Impact}: Chapter 5 (dynamic layer).
\end{proposal}
